% from ega0-20-fr.tex
\oldpage[0\textsubscript{IV-1}]{124}
\begin{definition}[20.4.2]
\label{0.20.4.2-fr}
On appelle \emph{$B$-algèbre augmentée des parties principales d'ordre $1$ de
$B$ par rapport à $A$} et l'on note $\mathrm{P}^1_{B/A}$ la $A$-algèbre
topologique quotient
\begin{equation}
\tag{20.4.2.1}
\label{0.20.4.2.1-fr}
  \mathrm{P}^1_{B/A}=(B\otimes_A B)/\mathfrak{J}^2
\end{equation}
munie de la structure de $B$-algèbre topologique définie par l'homomorphisme
$\bar{j}_1:B\to\mathrm{P}^1_{B/A}$ (déduit de $j_1$ par composition avec
l'homomorphisme canonique $B\otimes_A B\to\mathrm{P}^1_{B/A}$), et de
l'augmentation de $B$-algèbre $\varepsilon:\mathrm{P}^1_{B/A}\to B$ (aussi
notée $\varepsilon_{B/A}$) déduite de $p$ par passage au quotient.

Comme $p(b\otimes1)=b$ par définition, il est clair que $\varepsilon$ est
bien une augmentation de $B$-algèbre.
\end{definition}

\begin{definition}[20.4.3]
\label{0.20.4.3-fr}
Le noyau de l'augmentation $\varepsilon:\mathrm{P}^1_{B/A}\to B$,
\begin{equation}
\tag{20.4.3.1}
\label{0.20.4.3.1-fr}
  \Omega^1_{B/A}=\mathfrak{J}_{B/A}/(\mathfrak{J}_{B/A})^2
\end{equation}
muni de la topologie induite par celle de $\mathrm{P}^1_{B/A}$, qui en fait
un $B$-module topologique, est appelé le $B$-module des
\emph{$1$-différentielles} (ou simplement des \emph{différentielles}) de $B$
par rapport à $A$.
\end{definition}

On notera que la topologie de $\Omega^1_{B/A}$ est aussi la topologie quotient
de la topologie induite sur $\mathfrak{J}_{B/A}$ par celle de
$B\otimes_A B$ (Bourbaki, \emph{Top. gén.}, chap.~III, 3\textsuperscript{e}
éd., §~2, n\textsuperscript{o}~7, prop.~20). Si $B$ est discret il en est de
même de $\Omega^1_{B/A}$. On note $\widehat{\Omega}^1_{B/A}$ le
\emph{séparé complété} du $B$-module topologique $\Omega^1_{B/A}$.

Tout anneau topologique $B$ pouvant être considéré comme une
$\mathbf{Z}$-algèbre topologique (où $\mathbf{Z}$ est muni de la topologie
discrète), on peut définir le $B$-module topologique $\Omega^1_{B/\mathbf{Z}}$,
qu'on appelle parfois aussi le $B$-module des \emph{différentielles absolues}
sur $B$ et que l'on note $\Omega^1_B$. Si $B$ est une algèbre topologique sur
un corps premier $P$ (discret), on a
$B\otimes_P B=B\otimes_{\mathbf{Z}}B$, donc
$\Omega^1_{B/P}=\Omega^1_B$.

\begin{lemma}[20.4.4]
\label{0.20.4.4-fr}
Soient $A$ un anneau, $B$ une $A$-algèbre. L'idéal $\mathfrak{J}_{B/A}$ de
$B\otimes_A B$ est engendré par les éléments $1\otimes s-s\otimes1$, où $s$
parcourt un ensemble de générateurs de la $A$-algèbre $B$.
\end{lemma}

\begin{proof}
Il est clair que pour tout $x\in B$, on a
$x\otimes1-1\otimes x\in\mathfrak{J}$ ; d'autre part, quels que soient $x$,
$y$ dans $B$, on a
$x\otimes y=xy\otimes1+(x\otimes1)(1\otimes y-y\otimes1)$. Si
$\sum_i(x_i\otimes y_i)\in\mathfrak{J}$, on a par définition
$\sum_i x_i y_i=0$, donc
\begin{equation}
\tag{20.4.4.1}
\label{0.20.4.4.1-fr}
  \sum_i(x_i\otimes y_i)=
  \sum_i(x_i\otimes1)(1\otimes y_i-y_i\otimes1)
\end{equation}
ce qui prouve que $\mathfrak{J}$ est l'idéal engendré par les éléments
$1\otimes x-x\otimes1$. En outre, si $x=st$, on a
\begin{equation}
\tag{20.4.4.2}
\label{0.20.4.4.2-fr}
  x\otimes1-1\otimes x=(s\otimes1)(t\otimes1-1\otimes t)
  +(s\otimes1-1\otimes s)(1\otimes t)
\end{equation}
ce qui termine aussitôt la démonstration par récurrence.
\end{proof}

\begin{proposition}[20.4.5]
\label{0.20.4.5-fr}
Soient $A$ un anneau topologique, $B$ une $A$-algèbre topologique. La
topologie de $\Omega^1_{B/A}$ est moins fine que la topologie déduite de celle
de $B$ (19.0.2) ; si dans $B$ le carré de tout idéal ouvert est ouvert, ces
deux topologies sont identiques.
\end{proposition}

\begin{proof}
La première assertion est triviale, la topologie produit tensoriel sur
$B\otimes_A B$ étant moins fine que la topologie déduite de celle de $B$ ;
\emph{a fortiori} la topologie induite sur
$\mathfrak{J}=\mathfrak{J}_{B/A}$

